\subsection{CHAMBERS}

DEFINITION 2. A chamber of E relative to H (or simply a chamber if there is no ambiguity regarding H) is a facet of E relative to H that is not contained in any hyperplane belonging to H.

Let U be the open set in E consisting of the points that do not belong to any hyperplane of H; since a hyperplane of H must contain any facet that it meets, the chambers are the facets contained in U; every chamber is a convex (hence connected) open subset of E by Prop. 3, (i); since the chambers form a partition of U, they are exactly the connected components of U. Every convex subset A of U is connected, and thus contained in a chamber, which is unique if A is non-empty. It is clear that the chambers are the facets with support E, and Prop. 3, (iii) shows that every chamber is the interior of its closure. Finally, let C be a chamber and A a non-empty subset of C; formulas (2) and (3) imply that

\[
C = \bigcap_ {H \in \mathfrak {H}} D _ {H} (A) = D _ {\mathfrak {H}} (A), \quad \overline {{C}} = \bigcap_ {H \in \mathfrak {H}} \overline {{D _ {H} (A)}}\tag{6}
\]

since \(D_{\mathrm{H}}(\mathbf{A}) = D_{\mathrm{H}}(a)\) for all \(a \in A\) .

PROPOSITION 5. Let C be a non-empty subset of E. Assume that there exists a subset \(\mathfrak{H}'\) of \(\mathfrak{H}\) with the following properties:

\begin{enumerate}
\def\labelenumi{\alph{enumi})}
\item
  For any \(\mathrm{H} \in \mathfrak{H}'\) , there exists an open half-space \(\mathrm{D}_{\mathrm{H}}\) bounded by \(\mathrm{H}\) such that \(\mathrm{C} = \bigcap_{\mathrm{H} \in \mathfrak{H}'} \mathrm{D}_{\mathrm{H}}\) .
\item
  The set C does not meet any hyperplane belonging to \(\mathfrak{H} - \mathfrak{H}'\) .
\end{enumerate}

Under these conditions, C is a chamber defined by \(\mathfrak{H}\) in E, and \(\mathrm{D_H} = \mathrm{D_H(C)}\) for all \(\mathrm{H} \in \mathfrak{H}\) .

Properties a) and b) show that C is a convex subset of U; hence there is a chamber \(C'\) with \(C \subset C'\) . Since \(C \subset D_{H}\) , we have \(D_{H} = D_{H}(C)\) for all H in \(H'\) , hence \(C = D_{H'}(C) \supset D_{H}(C)\) since \(H' \subset H\) ; we have \(D_{H}(C) = C'\) by (6), hence \(C \supset C'\) . Finally therefore, \(C = C'\) .

PROPOSITION 6. Every point of E is in the closure of at least one chamber.

If E reduces to a single point, this is clear. Otherwise, let \(a \in E\) and let \(H_{1}, \ldots, H_{m}\) be the hyperplanes of H containing a. Since H is locally finite, there exists a neighbourhood V of a that does not meet any hyperplane of H other than \(H_{1}, \ldots, H_{m}\) . Let D be a straight line passing through a and not contained in any of the \(H_{i}\) ; if \(x \in D\) , \(x \neq a\) , and x is sufficiently close to a, the open segment \(|ax|\) is contained in V and does not meet any of the \(H_{i}\) . Then \(|ax| \subset U\) ; since \(|ax|\) is connected, it is contained in a chamber C, hence \(a \in \overline{C}\) .

PROPOSITION 7. Let L be an affine subspace of E and \(\Omega\) a non-empty open subset of L.

\begin{enumerate}
\def\labelenumi{(\roman{enumi})}
\item
  There exists a point \(a\) in \(\Omega\) that does not belong to any of the hyperplanes of \(\mathfrak{H}\) that do not contain \(L\) .
\item
  If \(\mathrm{L}\) is a hyperplane and \(\mathrm{L} \notin \mathfrak{H}\) , there exists a chamber that meets \(\Omega\) .
\item
  If \(\mathrm{L}\) is a hyperplane and \(\mathrm{L} \in \mathfrak{H}\) , there exists a point \(a\) in \(\Omega\) that does not belong to any hyperplane \(\mathrm{H} \neq \mathrm{L}\) of \(\mathfrak{H}\) .
\end{enumerate}

Denote by N the set of hyperplanes H with \(H \in H\) and \(L \notin H\) , and by L the set of hyperplanes of the affine space L of the form \(L \cap H\) with \(H \in N\) . It is clear that L is a locally finite set of hyperplanes in L, and Prop. 6 shows that \(\Omega\) meets a chamber \(\Gamma\) defined by L in L. If a is a point of \(\Gamma \cap \Omega\) , then \(a \notin H\) for all \(H \in N\) , hence (i).

Assume now that L is a hyperplane; any hyperplane containing L is then equal to it, so we may distinguish two cases:

\begin{enumerate}
\def\labelenumi{\alph{enumi})}
\item
  \(\mathbf{L} \notin \mathfrak{H}\) : then \(\mathfrak{N} = \mathfrak{H}\) , and we have \(a \notin \mathbf{H}\) for all \(\mathbf{H} \in \mathfrak{H}\) ; thus \(a\) belongs to a chamber defined by \(\mathfrak{H}\) in \(\mathbf{E}\) , hence (ii).
\item
  \(\mathrm{L} \in \mathfrak{H}\) : then \(\mathfrak{N} = \mathfrak{H} - \{\mathrm{L}\}\) , hence (iii).
\end{enumerate}

